% status: SKELETON
% Chapter 31 of the second edition. Plan: docs/STRUCTURE.md. Rules: AUTHORING.md.
\chapter{Tensor, Pipeline and Context Parallelism}\label{ch:model-parallelism}

\begin{ChapterIntent}
When a model does not fit on one device, or one device is too slow, the
model must be split. Tensor parallelism splits each matrix and pays collective
communication in every layer; pipeline parallelism splits the layers and pays
in bubbles; context parallelism splits the sequence for long prefill. This
chapter derives each, prices its communication in bytes against the
interconnect's bandwidth, and shows how to choose a layout for a given model,
context and latency target.
\end{ChapterIntent}

\OpeningSection{Two all-reduces per layer}

\NapkinSection{Napkin math: collective bytes against interconnect bandwidth}

\section{Why split a model}

\section{Tensor parallelism}

\section{The collectives each layer pays}

\section{Pipeline parallelism and its bubbles}

\section{Context parallelism and ring attention}

\section{Choosing a layout}

\section{Interconnects}

\LedgerSection

\LabSection{simulate tensor parallelism with threads}

\HappenedSection{partitioning PaLM for inference}

\FrontierSection{}

\ExercisesSection
